% !TEX root = ../../main.tex
% C12.2 - Limitations. 12.2.1 = Vietnamese 9.2 translated (Part I); 12.2.2 = research limitations (Part II), verbatim.
\section{Limitations}
\label{sec:9.2}

\subsection{Part I: The Application}
\label{sec:12.2.1}

The experiments confirm the limitations stated in Section~\ref{sec:1.4}. The largest is that search by text and by attributes is not yet effective on the test data: no query has a correct result in the top 16 (Section~\ref{sec:8.4}). The cause is the domain gap between RaSa's training data and person images cropped automatically from cameras.

On a CPU-only machine, processing is about seven times slower than real time. The system has to process the cameras in turn, part of the frames of a live stream are skipped, and the demonstration data had to be indexed from video files (Section~\ref{sec:7.2}). RTSP data is also only encoded and written at the end of each session, rather than as soon as each appearance ends, so there is a delay between a person appearing in front of a camera and becoming searchable (Section~\ref{sec:5.2.1}). RTSP frames are also timestamped when read rather than by the stream, so slow processing widens the gap between processed frames, ending tracks early and splitting appearances (Section~\ref{sec:3.3.5}); video files are not affected. The first search needs a warm-up to load the encoder, and memory is tight: with the server and the pipeline each loading their own models, processing uses about 8~GiB, nearly exhausting a 16~GB machine with other applications open (Section~\ref{sec:8.7}).

Finally, the evaluation is incomplete: the quality of search by image and by text at the current confidence threshold has not been measured again (Section~\ref{sec:8.4}). BoT-SORT is integrated but has not been used for indexing or evaluated.

\subsection{Part II: The Reasoning Layer}
\label{sec:12.2.2}

% !TEX root = ../../main.tex
% Part II (Nguyen Huu Cuong) - merged from the research report (report_latex_v3), see MERGE_NOTES.md.
% Style pass 2026-09-30: em dashes and rhetorical sentences removed; technical content, numbers and references unchanged.
% Used by Chapter 12, section 12.2 (Part II limitations).
\label{sec:r-limits}
\paragraph{Model.} The reasoning LLM is Flash-Lite 3.5, chosen for its larger
free-tier allowance rather than its quality. Section~\ref{sec:r-rerun} confirms
that this is a limitation: it recovers four of five ranking failures by switching
the two ranking calls to a larger model and changing nothing else. A full sweep on
the larger model is the most direct improvement available and has not been run;
the re-run covers five oracle-selected queries, not the sample.

\paragraph{Data.} Token limits restricted the evaluation to 100 of the 496
uniquely identifiable identities, so practical usability is not yet established,
and the subset itself is an easier population than the full 1501 identities
(Section~\ref{sec:r-dataset}). At $n=100$ a percentage point is one identity, and
the R@1 ordering in Table~\ref{tab:r-main} rests on differences of 3 and 8
identities. Extending the sweep requires no new work: the work-queue notebook
already selects the next identities to run, giving priority to queries whose
ground truth is inside the pool but outside CLIP's top 10.

\paragraph{The blocks this part did not change.} Detection and tracking were filled
with established defaults and left unchanged (Section~\ref{sec:r-choices}), and
their failures still reach the search block: a person the detector missed has no
crop, and a track merged by the tracker contains two people under one identity.
This part measures neither. A tracker-error analysis over the same 100 identities
is the next measurement to make, and it would move part of what
Section~\ref{sec:r-error} counts as a reasoning residue into an upstream category.

\paragraph{The refine round.} Four stages beat two at R@10 and lose to two at R@1
and R@5, and three of the five losses in Table~\ref{tab:r-lostwhy} were caused by
the refine round. Whether stages 3 and 4 should return a full ranking, or only
re-order the candidates the coarse round was unsure about, is still open.

\paragraph{Rarity versus mismatch cost.} The rarity table tells the ranker what a
query value is worth when it separates candidates. It does not yet state the
converse: a \emph{binary} attribute such as gender is ``almost nothing'' by rarity
(both values are common), yet a gender mismatch is the strongest exclusion
available. On one measured query, 40 of 40 reasons led with the lower colour and
gender appeared in 2, while gender split the pool 44 to 31. This asymmetry between
rarity and mismatch cost remains to be handled.

\paragraph{Output truncation.} At 25 candidates per 2A call, replies occasionally
stop at the output limit in the middle of a unit. The limit is currently detected
after the fact from the API's finish reason rather than requested explicitly, and
the two direct fixes, sending an explicit output cap or reducing the batch to 15
candidates, have not yet been compared.

\paragraph{Single phrasing.} Each identity is queried with one descriptive
sentence. Robustness to paraphrase is not measured, although the pipeline already
supports several phrasings per identity.

